\chapter{Etude préalable}

\section{Description de l'existant}
L'analyse initiale a montré une gestion fragmentée des données pédagogiques et administratives. Les informations relatives aux absences, demandes, réclamations et notifications n'étaient pas centralisées dans une même architecture de services.

L'existant présentait également un manque de standardisation dans l'exposition des fonctionnalités côté API.

\section{Critiques de l'existant}
Les limites observées sont les suivantes :
\begin{itemize}
	\item redondance des informations entre plusieurs modules ;
	\item lenteur dans le traitement des opérations administratives ;
	\item difficulté de traçabilité des actions ;
	\item intégration frontend/backend moins fluide en l'absence de contrats API homogènes.
\end{itemize}

\section{Solutions proposées}
Pour répondre à ces insuffisances, la solution proposée repose sur :
\begin{itemize}
	\item une architecture en couches \textit{Controller/Service/Repository} ;
	\item un modèle de données unifié basé sur des entités JPA ;
	\item des endpoints REST standardisés ;
	\item une gestion globale des erreurs via un gestionnaire d'exceptions centralisé.
\end{itemize}

\section{Développement de la solution finale}
La solution finale backend intègre les modules critiques suivants :
\begin{itemize}
	\item gestion des utilisateurs (admin, enseignant, étudiant) ;
	\item gestion des départements, cours et structures pédagogiques ;
	\item suivi des absences et demandes ;
	\item traitement des réclamations ;
	\item système de notifications ;
	\item services de persistance et de gestion des pièces jointes.
\end{itemize}
